%%%%%%%%%%%%%%%%%%%%%%%%%%%%%%%%%%%%%%%%%%%%%%%%%%%%%%%%%%%%%%%%%%%%%%%%%%%%%%%%
% paper_en.tex
% Nature-style research article format (A4)
% Author: Xiaohui Li
%%%%%%%%%%%%%%%%%%%%%%%%%%%%%%%%%%%%%%%%%%%%%%%%%%%%%%%%%%%%%%%%%%%%%%%%%%%%%%%%

\documentclass[12pt,a4paper]{article}

% Page geometry
\usepackage[a4paper,left=25mm,right=25mm,top=25mm,bottom=25mm]{geometry}

% Encoding and fonts
\usepackage[utf8]{inputenc}
\usepackage[T1]{fontenc}
\usepackage{lmodern}

% Math and tables
\usepackage{amsmath,amssymb}
\usepackage{booktabs}
\usepackage{array}

% Figures
\usepackage{graphicx}

% Hyperlinks
\usepackage[colorlinks=true,linkcolor=blue,citecolor=blue,urlcolor=blue]{hyperref}

% Line spacing
\usepackage{setspace}
\onehalfspacing

% Title formatting
\usepackage{titlesec}
\titleformat{\section}{\large\bfseries}{\thesection.}{0.5em}{}
\titleformat{\subsection}{\normalsize\bfseries}{\thesubsection}{0.5em}{}

%%%%%%%%%%%%%%%%%%%%%%%%%%%%%%%%%%%%%%%%%%%%%%%%%%%%%%%%%%%%%%%%%%%%%%%%%%%%%%%%
\title{
  \textbf{Quantum-State Homophily as a Candidate Origin of Gravitation:}\\
  \textbf{A Hypothesis, Formal Framework, and Testable Predictions}
}

\author{Xiaohui Li\\
Independent Researcher, Baoji 721000, China\\
Correspondence Xiaohui Li: \texttt{danceluck@163.com}\\
Xiaohui Li ORCID:0009-0003-8233-7191
}

\date{}

%%%%%%%%%%%%%%%%%%%%%%%%%%%%%%%%%%%%%%%%%%%%%%%%%%%%%%%%%%%%%%%%%%%%%%%%%%%%%%%

\begin{document}

\maketitle
\thispagestyle{fancy}

%%%%%%%%%%%%%%%%%%%%%%%%%%%%%%%%%%%%%%%%%%%%%%%%%%%%%%%%%%%%%%%%%%%%%%%%%%%%%%%

\begin{abstract}
\noindent

Although general relativity describes gravitation with exquisite precision, its reconciliation with the quantum realm remains unresolved. Here we propose a candidate organizing principle---the Quantum-State Phase Attraction--Repulsion (QSPAR) hypothesis---wherein identical quantum states attract and complementary (orthogonal) quantum states repel. We formalize the concept of a ``material line'' as a coherent bundle of in-phase quanta and show that, in the unpolarized mixed-state limit, the model is consistent with Newtonian gravitation, whereas ordered, same-phase material lines can generate a repulsive correction that we identify with a tunable anti-gravity effect. The framework yields three falsifiable predictions: (i) a chiral weight differential between enantiomers, (ii) a measurable reduction of effective weight for rapidly rotating polarized matter, and (iii) orbit-period anomalies for two-body celestial systems. We outline numerical benchmarks against Cavendish-type, equivalence-principle and inverse-square-law tests, emphasizing empirical falsifiability.
\end{abstract}

\vspace{0.5em}
\noindent\textbf{Keywords:} quantum state; material line; gravitation; anti-gravity; chirality; falsifiable prediction; coarse-graining

\section*{Main Text}

\subsection*{1. Introduction}

Gravitation was the first interaction to be described systematically~\cite{newton1687}, and Einstein's general theory of relativity geometrized it as spacetime curvature, passing extremely stringent tests in the solar system~\cite{einstein1916,will2014}. Yet the quantum origin of gravity---how gravitation emerges from quantum states---remains one of the deepest open problems in theoretical physics~\cite{feynman1995}. Recent advances in quantum coherence, entanglement, and macroscopic order parameters offer a fresh vocabulary with which to reconsider gravity~\cite{zurek2003}.

Inspired by the behaviour of magnetic field lines~\cite{parker1979} and by homophily in social networks~\cite{mcpherson2001}, we propose the Quantum-State Phase Attraction--Repulsion (QSPAR) hypothesis: identical quantum states attract, whereas complementary quantum states repel. We elevate this intuitive picture into a formal ``material-line'' model and show that (i) in ordinary unpolarized matter the attractive channel dominates, reproducing Newtonian gravity, and (ii) by helically ordering mixed quantum states into coherent material lines, a repulsive channel can be opened, enabling anti-gravity engineering. The value of QSPAR lies not in replacing general relativity, but in providing a conceptually transparent, quantitatively testable framework rooted in quantum-state structure.

\subsection*{2. Results}

\subsubsection*{2.1 Formal QSPAR framework}

Let a coarse-grained quantum state be labelled by a phase $\phi \in [0, 2\pi)$. A ``material line'' $L(\phi)$ is defined as a coherent bundle of quanta sharing the same phase $\phi$. The interaction potential between two material lines $a$ and $b$ separated by distance $r$ is written as
\begin{equation}
U_{ab}(r) = -\frac{G m_a m_b}{r}\cos^2\!\left(\frac{\phi_a-\phi_b}{2}\right)
            +\frac{K(\omega_{ab}) m_a m_b}{r}\sin^2\!\left(\frac{\phi_a-\phi_b}{2}\right)e^{-r/\lambda},
\label{eq:qspar}
\end{equation}
where $G$ is an effective gravitational constant, $\lambda$ is a characteristic cutoff length of the repulsive interaction, and $K(\omega_{ab})$ is a repulsive coupling that depends on the relative angular velocity $\omega_{ab}$. This potential satisfies two limits:
\begin{itemize}
  \item when $\phi_a=\phi_b$ and $\omega_{ab}=0$, $U_{ab}=-Gm_am_b/r$, giving pure attraction;
  \item when $|\phi_a-\phi_b|=\pi$ and relative motion is present, the second term is activated, giving repulsion.
\end{itemize}
This expression mathematizes the ``same phase attracts, opposite phase repels'' intuition and naturally introduces a velocity-dependent anti-gravity channel.

\subsubsection*{2.2 Mixed states, polarization, and the Newtonian limit}

In ordinary matter, quantum phases are typically disordered. We define a polarization order parameter $P\in[0,1]$ that quantifies the degree of material-line ordering: $P=0$ corresponds to a fully mixed state and $P=1$ to a fully ordered array of material lines. Ensemble-averaging over randomly oriented material lines causes the $\sin^2$ term to cancel, while the $\cos^2$ term averages to $1/2$, yielding an effective gravitational potential
\begin{equation}
\langle U\rangle = -\frac{G_{\rm eff}m_a m_b}{r}, \qquad G_{\rm eff}\approx \frac{G}{2}.
\end{equation}
By re-identifying $G_{\rm eff}$ with the observed Newtonian constant $G_N$, the model becomes indistinguishable from Newtonian gravity for static, unpolarized matter. This limit guarantees consistency with Cavendish experiments and planetary dynamics.

\subsubsection*{2.3 Material-line mechanism of anti-gravity}

If the mixed quantum state of a body is ``spiralled upward'' into ordered material lines---for example by strong magnetic fields, rapid rotation, or chiral-crystal engineering---then $P>0$ and the repulsive channel opens. For a polarized disk of mass $m$ rotating with angular velocity $\omega$ relative to Earth's gravitational field, the fractional weight correction is approximately
\begin{equation}
\frac{\Delta W}{W} \approx \beta\left(\frac{\omega}{\omega_0}\right)^2 P^2,
\end{equation}
where $\beta$ is a material-dependent dimensionless coupling and $\omega_0$ is a reference angular frequency. This predicts that, for sufficiently large $P$ and $\omega$, the disk exhibits a measurable weight reduction---the conceptual core of an anti-gravity device.

\subsubsection*{2.4 Falsifiable predictions}

QSPAR leads to three directly testable predictions:

\noindent\textbf{Prediction 1 (chiral weight differential).} Biology is dominated by left-handed chiral molecules. If chirality is a macroscopic manifestation of quantum-state phase, left- and right-handed enantiomers should exhibit a tiny weight difference, predicted to be $\Delta g/g < 10^{-12}$, within reach of modern torsion-balance experiments~\cite{schlamminger2008}.

\noindent\textbf{Prediction 2 (weight reduction of a rotating polarized disk).} A disk made of highly polarized magnetic or chiral material and rotated in the range $10^4$--$10^5$~rad s$^{-1}$ should display an effective weight change at the $10^{-6}$--$10^{-4}$ level. Such an experiment can be performed in a vacuum torsion balance or a cryogenic levitation platform.

\noindent\textbf{Prediction 3 (celestial orbital anomaly).} If the Earth--Moon system can be approximated as two complementary but relatively moving material lines, their potential energy will contain a small repulsive correction, producing a long-term drift in the lunar orbital period relative to the Keplerian value. Lunar laser-ranging data~\cite{will2014} limit this drift to the centimetre level, thereby constraining $\lambda$ and $K$.

\subsubsection*{2.5 Benchmarking against established gravity tests}

To ensure that QSPAR does not conflict with established gravitational physics, we compare it with classical benchmark tests (Table~\ref{tab:benchmarks}).

\begin{table}[htbp]
\centering
\caption{Benchmarking QSPAR against established gravity tests.}
\label{tab:benchmarks}
\begin{tabular}{@{}p{4.2cm}p{4.5cm}p{5.8cm}@{}}
\toprule
Experiment/Observation & Current constraint & Implication for QSPAR \\
\midrule
Inverse-square test (cm) & $|\delta\alpha| < 10^{-3}$~\cite{adelberger2003} & Repulsive cutoff $\lambda \lesssim 1$~cm or $K$ strongly suppressed \\
Equivalence principle & $\eta < 10^{-15}$~\cite{schlamminger2008} & Chiral/phase coupling must lie below this bound \\
Lunar laser ranging & Distance drift $<$ few cm over decades~\cite{will2014} & Long-range correction negligible or constrained \\
Sub-mm inverse-square & $|\alpha| < 10^{-4}$ at 100~$\mu$m~\cite{hoyle2004} & $\lambda \lesssim 100~\mu$m or repulsion is purely local \\
\bottomrule
\end{tabular}
\end{table}

\subsection*{3. Discussion}

QSPAR offers a conceptual path toward understanding gravity from quantum-state phase. It does not deny the macroscopic success of general relativity; rather, it attempts to provide a microscopic mechanism for the gravitational constant $G$: $G$ may reflect the high degree of quantum-state homogeneity of cosmic matter. The framework's most attractive feature is falsifiability---if future experiments on chiral samples, rotating polarized bodies, or celestial orbits fail to reveal the predicted effects, the hypothesis can be ruled out.

Nevertheless, QSPAR remains a phenomenological model. It has not yet been derived from first principles and lacks a clear connection to quantum field theory and gauge symmetry; the functions $K(\omega)$ and $\lambda$ in the potential still require microscopic or experimental calibration. Moreover, the analogies in the original manuscript concerning nations, languages, and interpersonal relations~\cite{mcpherson2001} should be treated only as heuristic metaphors for coarse-grained quantum states in social systems, not as physical evidence.

Future work should include: (1) constructing a microscopic Lagrangian based on quantum field theory; (2) designing high-precision chiral-weight experiments; (3) testing the hypothesis that ``the universe as a whole occupies a single quantum state'' in cosmological structure-formation simulations; and (4) exploring strongly correlated systems such as topological insulators and spin liquids for macroscopic quantum coherence that might be engineered into material lines.

\subsection*{4. Methods}

\subsubsection*{4.1 Coarse-grained lattice simulation}

To test QSPAR quantitatively, we propose a reproducible lattice-simulation protocol:

\begin{enumerate}
  \item Discretize space into a cubic lattice of spacing $a$; assign each site a mass $m_i$ and phase $\phi_i$.
  \item For each pair $(i,j)$, compute the force from the potential: $\mathbf{F}_{ij}=-\nabla U_{ij}(r_{ij})$.
  \item Evolve the system with a velocity-Verlet integrator, using a time step $\Delta t$ satisfying the stability condition $\Delta t < \sqrt{a^3/(Gm)}$.
  \item Compute the total potential energy, effective gravitational constant, and polarization order parameter $P$.
\end{enumerate}

\noindent\textit{Pseudocode} (Python/NumPy sketch):
\begin{verbatim}
import numpy as np

def qspar_force(r, phi_i, phi_j, m_i, m_j,
                G, K0, lam, omega, omega0):
    dphi = phi_i - phi_j
    dist = np.linalg.norm(r)
    attr = -G*m_i*m_j/dist**2 * np.cos(dphi/2)**2
    rep  =  K0*(omega/omega0)**2 * m_i*m_j/dist**2 \
            * np.sin(dphi/2)**2 * np.exp(-dist/lam)
    return (attr + rep) * r/dist
\end{verbatim}

\subsubsection*{4.2 Parameter calibration}

Table~\ref{tab:params} lists the reference parameters for an Earth--Moon simulation. All constants are taken from CODATA 2018 or the NASA JPL ephemerides.

\begin{table}[htbp]
\centering
\caption{Reference parameters for Earth--Moon QSPAR simulations.}
\label{tab:params}
\begin{tabular}{@{}llll@{}}
\toprule
Parameter & Symbol & Value & Source \\
\midrule
Earth mass & $M_E$ & $5.972\times10^{24}$~kg & CODATA \\
Moon mass & $M_M$ & $7.348\times10^{22}$~kg & CODATA \\
Earth--Moon distance & $r$ & $3.844\times10^{8}$~m & NASA JPL \\
Gravitational constant & $G$ & $6.67430\times10^{-11}$~m$^3$ kg$^{-1}$ s$^{-2}$ & CODATA \\
Reference angular frequency & $\omega_0$ & $10^4$~rad s$^{-1}$ & Experimental design \\
Repulsive cutoff & $\lambda$ & $10^{-2}$~m (tunable) & Assumption \\
\bottomrule
\end{tabular}
\end{table}

\subsubsection*{4.3 Validation strategy}

We first reproduce a two-body Keplerian orbit in the fully mixed state ($P=0$), requiring agreement with Newtonian theory at the $10^{-6}$ level. We then introduce a small polarization $P=10^{-3}$ and measure any long-term orbital drift, comparing it with lunar laser-ranging data. This two-step validation ensures correctness in the known limit and provides a quantifiable window for new physics.

\section*{Data Availability}

No new experimental data were generated in this study. The simulation parameters in Table~\ref{tab:params} are derived from public databases: CODATA 2018 and the NASA JPL ephemerides. The conceptual ``material-line'' simulation parameter set is provided as Supplementary Data~1.



%%%%%%%%%%%%%%%%%%%%%%%%%%%%%%%%%%%%%%%%%%%%%%%%%%%%%%%%%%%%%%%%%%%%%%%%%%%%%%%%
% References
%%%%%%%%%%%%%%%%%%%%%%%%%%%%%%%%%%%%%%%%%%%%%%%%%%%%%%%%%%%%%%%%%%%%%%%%%%%%%%%%
\begin{thebibliography}{10}

\bibitem{newton1687}
Newton I.
\newblock {\em Philosophi{\ae} Naturalis Principia Mathematica}.
\newblock London: Royal Society; 1687.

\bibitem{einstein1916}
Einstein A.
\newblock Die Grundlage der allgemeinen Relativit\"{a}tstheorie.
\newblock {\em Annalen der Physik}. 1916;354(7):769--822.

\bibitem{feynman1995}
Feynman RP, Morinigo FB, Wagner WG.
\newblock {\em Feynman Lectures on Gravitation}.
\newblock Reading, MA: Addison-Wesley; 1995.

\bibitem{will2014}
Will CM.
\newblock The confrontation between general relativity and experiment.
\newblock {\em Living Reviews in Relativity}. 2014;17:4.

\bibitem{adelberger2003}
Adelberger EG, Heckel BR, Nelson AE.
\newblock Tests of the gravitational inverse-square law below the dark-energy length scale.
\newblock {\em Annual Review of Nuclear and Particle Science}. 2003;53:77--121.

\bibitem{zurek2003}
Zurek WH.
\newblock Decoherence, einselection, and the quantum origins of the classical.
\newblock {\em Reviews of Modern Physics}. 2003;75(3):715--775.

\bibitem{parker1979}
Parker EN.
\newblock {\em Cosmical Magnetic Fields: Their Origin and Their Activity}.
\newblock Oxford: Clarendon Press; 1979.

\bibitem{mcpherson2001}
McPherson M, Smith-Lovin L, Cook JM.
\newblock Birds of a feather: homophily in social networks.
\newblock {\em Annual Review of Sociology}. 2001;27:415--444.

\bibitem{hoyle2004}
Hoyle CD, Schmidt U, Heckel BR, et~al.
\newblock Submillimeter tests of the gravitational inverse-square law.
\newblock {\em Physical Review D}. 2004;70:042004.

\bibitem{schlamminger2008}
Schlamminger S, Choi KY, Wagner TA, Gundlach JH, Adelberger EG.
\newblock Test of the equivalence principle using a rotating torsion balance.
\newblock {\em Physical Review Letters}. 2008;100(4):041101.

\end{thebibliography}
\textbf{\textbf{Declarations}}

  \textbf{\textbf{Acknowledgements}}

The author thanks colleagues for discussions on quantum aggregates, consciousness, emergence, and dimensionality. 

 \textbf{\textbf{Funding Statement}}:This research did not receive any specific grant from funding agencies in the public,commercial,or not-for-profit sectors.

Currently, there are no sponsors. Or I haven't found any sponsors or organizations yet. However, if I can publish it, I will cover part of the costs myself. Or it will be easier to find sponsors. (For example, the German company E.ON. I have previously found international buyers for their wind turbines.)

 \textbf{\textbf{Competing Interests}}

The author declares no competing interests.

 \textbf{\textbf{Author Contributions}}

Xiaohui Li: Conceptualization, Methodology, Investigation, Formal analysis, Writing – Original Draft, Visualization, Writing – Review \& Editing.

AI assisted me in rewriting the original text into an article that meets the submission requirements of Journal of the American Mathematical Society. No content was added or omitted.

 The author is solely responsible for the work described in this manuscript and has approved the final version for publication.

 \textbf{\textbf{Clinical trial:}}not applicable.

 \textbf{\textbf{Ethics and Consent to Participate declarations}}: not applicable.

 Hi,

This is an original work by the author. There is only one author:Xiaohui Li.

The article accurately reflects the scientific results.

Xiaohui Li agreed to submit the manuscript.

This manuscript has not been published in any other journals.

The author accepts the established procedure for selecting and publishing manuscripts.

good

Xiaohui Li



 \textbf{\textbf{Clinical trial number}}: not applicable.
\end{document}
%%%%%%%%%%%%%%%%%%%%%%%%%%%%%%%%%%%%%%%%%%%%%%%%%%%%%%%%%%%%%%%%%%%%%%%%%%%%%%%%
